% --- Archon physics formalization source begin ---
% archon:physics
% archon:covers PhyXMiniProblems/problem_phyx_mini_0564.lean
% archon:source-report reports/phyx_mini/problem_phyx_mini_0564.source.json

\chapter{Physics problem phyx\_mini\_0564}
\label{ch:PhyXMiniProblems_problem_phyx_mini_0564}

\paragraph{Problem source.}
\#\# Physical scenario

An early method testing Maxwell's theoretical prediction for the distribution of molecular speeds is shown in figure. In 1925 Otto Stern used a beam of \textbackslash{}(\textbackslash{}text\{Bi\}\_2\textbackslash{}) molecules emitted from an oven at 850 K. The beam defined by slit \textbackslash{}(\textbackslash{}text\{S\}\_1\textbackslash{}) was admitted into the interior of a rotating drum via slit \textbackslash{}(\textbackslash{}text\{S\}\_2\textbackslash{}) in the drum wall. The identical bunches of molecules thus formed struck and adhered to a curved glass plate fixed to the interior drum wall, the fastest molecules striking near A, which was opposite \textbackslash{}(\textbackslash{}text\{S\}\_2\textbackslash{}), the slowest near B, and the others in between depending on their speeds. The density of the molecular deposits along the glass plate was measured with a densitometer. The density (proportional to the number of molecules) plotted against distance along the glass plate (dependent on \textbackslash{}( v \textbackslash{})) made possible determination of the speed distribution. If the drum is 10 cm in diameter and is rotating at 6250 rpm.

\#\# Question

Find the distance from A where molecules traveling at \textbackslash{}( \textbackslash{}langle v \textbackslash{}rangle \textbackslash{}) will strike.

\#\# Answer choices

- A: A: \textbackslash{}( 0.04625 \textbackslash{}text\{rev\} \textbackslash{})

- B: B: \textbackslash{}( 0.05667 \textbackslash{}text\{rev\} \textbackslash{})

- C: C: \textbackslash{}( 0.05317 \textbackslash{}text\{rev\} \textbackslash{})

- D: D: \textbackslash{}( 0.05021 \textbackslash{}text\{rev\} \textbackslash{})

\#\# Figure caption (auxiliary; use the image as primary evidence)

The image features a diagram illustrating an experimental setup.

- **Oven:** The setup begins with an oven on the right side which seems to emit or influence some form of beam.
  
- **Beam Path:** There is a linear path extending from the oven passing through two slits, S1 and S2.

- **Slits (S1 and S2):** Two vertical slits are spaced along the beam's path, acting as barriers with openings.

- **Drum and Glass Plate:** After passing through the slits, the path enters a circular drum labeled as “Drum.” Inside the drum, there is a “Glass plate” rotating about its center with an angular velocity labeled as \textbackslash{}( w \textbackslash{}).

- **Quadrants (A and B):** The drum is marked with two specific points, A and B, likely indicating key locations on the glass plate. 

The image appears to illustrate an experimental arrangement involving interference or spectroscopy, with precise mechanical components aligning the pathway of the beam through strategic openings and onto rotating measuring equipment.

\#\# Dataset metadata

Quantum Phenomena; Physical Model Grounding Reasoning, Spatial Relation Reasoning

\paragraph{Recorded answer/context.}
D: D: \textbackslash{}( 0.05021 \textbackslash{}text\{rev\} \textbackslash{})

\paragraph{Figure/image path.}
/root/proposal\_for\_physic/science-mango/phyx\_mini\_run/phyx\_data/test\_image/564.png

\paragraph{Formalization target.}
create a compiling Lean file with sorry bodies at `PhyXMiniProblems/problem\_phyx\_mini\_0564.lean`. The Lean declarations must preserve the physical quantities, dimensions or dimensional roles, figure labels, governing-law hypotheses, and final relation expressed by this problem.
Use Mathlib/Physlib names found through LeanExplore where available. If a physics API is missing, introduce faithful local abstractions rather than scalar placeholder aliases.

\begin{theorem}[Physics formalization target]
\label{thm:physics:phyx_mini_0564:target}
\lean{PhyXMiniProblems.ProblemPhyXMini0564.problem_phyx_mini_0564}
\uses{def:physics:phyx-mini-0564:phyxminiproblems-problemphyxmini0564-lengthreadout, def:physics:phyx-mini-0564:phyxminiproblems-problemphyxmini0564-speedinmeterspersecond, def:physics:phyx-mini-0564:phyxminiproblems-problemphyxmini0564-rotationfrequencyreadout, def:physics:phyx-mini-0564:phyxminiproblems-problemphyxmini0564-sternspeeddistributionsetup, def:physics:phyx-mini-0564:phyxminiproblems-problemphyxmini0564-meanspeedstrikeoffsetrevolutions, def:physics:phyx-mini-0564:phyxminiproblems-problemphyxmini0564-matchessternexperimentscenario, def:physics:phyx-mini-0564:phyxminiproblems-problemphyxmini0564-matchessternproblemreadouts, def:physics:phyx-mini-0564:phyxminiproblems-problemphyxmini0564-usesbismuthdimerreferencedata, def:physics:phyx-mini-0564:phyxminiproblems-problemphyxmini0564-matchesprimarysternfigure, def:physics:phyx-mini-0564:phyxminiproblems-problemphyxmini0564-hasphysicalsternparameters, def:physics:phyx-mini-0564:phyxminiproblems-problemphyxmini0564-satisfiesmaxwellmeanspeedlaw, def:physics:phyx-mini-0564:phyxminiproblems-problemphyxmini0564-satisfiesrotatingdrumkinematics, def:physics:phyx-mini-0564:phyxminiproblems-problemphyxmini0564-satisfiesinversespeedsorting, def:physics:phyx-mini-0564:phyxminiproblems-problemphyxmini0564-satisfiesdensitometerlaw, def:physics:phyx-mini-0564:phyxminiproblems-problemphyxmini0564-displayedoffsetrevolutions, def:physics:phyx-mini-0564:phyxminiproblems-problemphyxmini0564-agreeswithinrevolutions, def:physics:phyx-mini-0564:phyxminiproblems-problemphyxmini0564-isuniquenearestdisplayedoffset}
The assigned autoformalize agent should translate this physics problem into Lean declarations in the covered file, with theorem and lemma proof bodies written as `by sorry`.
\end{theorem}
\begin{proof}
This is an autoformalization task, not a proof task. Produce statements that can later be proved by the physics prover without weakening the physical model.
\end{proof}

% --- Archon declaration topology begin ---
\section{Declaration topology}
\begin{definition}[Length Quantity]
  \label{def:physics:phyx-mini-0564:phyxminiproblems-problemphyxmini0564-lengthquantity}
  \lean{PhyXMiniProblems.ProblemPhyXMini0564.LengthQuantity}
  A nonnegative physical length independent of the unit used to read it
\end{definition}
\begin{proof}
  This is the declared model definition of Length Quantity.
\end{proof}
\begin{definition}[Time Quantity]
  \label{def:physics:phyx-mini-0564:phyxminiproblems-problemphyxmini0564-timequantity}
  \lean{PhyXMiniProblems.ProblemPhyXMini0564.TimeQuantity}
  A nonnegative physical duration
\end{definition}
\begin{proof}
  This is the declared model definition of Time Quantity.
\end{proof}
\begin{definition}[Temperature Quantity]
  \label{def:physics:phyx-mini-0564:phyxminiproblems-problemphyxmini0564-temperaturequantity}
  \lean{PhyXMiniProblems.ProblemPhyXMini0564.TemperatureQuantity}
  A nonnegative absolute temperature
\end{definition}
\begin{proof}
  This is the declared model definition of Temperature Quantity.
\end{proof}
\begin{definition}[Molar Mass Quantity]
  \label{def:physics:phyx-mini-0564:phyxminiproblems-problemphyxmini0564-molarmassquantity}
  \lean{PhyXMiniProblems.ProblemPhyXMini0564.MolarMassQuantity}
  A nonnegative molar mass. Amount of substance is treated as a fixed one-mole reference because Physlib's current Dimension has no mole coordinate
\end{definition}
\begin{proof}
  This is the declared model definition of Molar Mass Quantity.
\end{proof}
\begin{definition}[Rotation Frequency Quantity]
  \label{def:physics:phyx-mini-0564:phyxminiproblems-problemphyxmini0564-rotationfrequencyquantity}
  \lean{PhyXMiniProblems.ProblemPhyXMini0564.RotationFrequencyQuantity}
  A nonnegative rotation frequency, with revolutions counted dimensionlessly
\end{definition}
\begin{proof}
  This is the declared model definition of Rotation Frequency Quantity.
\end{proof}
\begin{definition}[Molar Gas Constant Quantity]
  \label{def:physics:phyx-mini-0564:phyxminiproblems-problemphyxmini0564-molargasconstantquantity}
  \lean{PhyXMiniProblems.ProblemPhyXMini0564.MolarGasConstantQuantity}
  The universal molar gas constant, with the mole count treated dimensionlessly: mass * length² * time⁻² * temperature⁻¹
\end{definition}
\begin{proof}
  This is the declared model definition of Molar Gas Constant Quantity.
\end{proof}
\begin{definition}[length Readout]
  \label{def:physics:phyx-mini-0564:phyxminiproblems-problemphyxmini0564-lengthreadout}
  \lean{PhyXMiniProblems.ProblemPhyXMini0564.lengthReadout}
  \uses{def:physics:phyx-mini-0564:phyxminiproblems-problemphyxmini0564-lengthquantity}
  Read a physical length in a selected length unit
\end{definition}
\begin{proof}
  This is the declared model definition of length Readout.
\end{proof}
\begin{definition}[time Readout]
  \label{def:physics:phyx-mini-0564:phyxminiproblems-problemphyxmini0564-timereadout}
  \lean{PhyXMiniProblems.ProblemPhyXMini0564.timeReadout}
  \uses{def:physics:phyx-mini-0564:phyxminiproblems-problemphyxmini0564-timequantity}
  Read a physical duration in a selected time unit
\end{definition}
\begin{proof}
  This is the declared model definition of time Readout.
\end{proof}
\begin{definition}[temperature Readout]
  \label{def:physics:phyx-mini-0564:phyxminiproblems-problemphyxmini0564-temperaturereadout}
  \lean{PhyXMiniProblems.ProblemPhyXMini0564.temperatureReadout}
  \uses{def:physics:phyx-mini-0564:phyxminiproblems-problemphyxmini0564-temperaturequantity}
  Read an absolute temperature in a selected temperature unit
\end{definition}
\begin{proof}
  This is the declared model definition of temperature Readout.
\end{proof}
\begin{definition}[molar Mass Readout]
  \label{def:physics:phyx-mini-0564:phyxminiproblems-problemphyxmini0564-molarmassreadout}
  \lean{PhyXMiniProblems.ProblemPhyXMini0564.molarMassReadout}
  \uses{def:physics:phyx-mini-0564:phyxminiproblems-problemphyxmini0564-molarmassquantity}
  Read molar mass in the selected mass unit per mole
\end{definition}
\begin{proof}
  This is the declared model definition of molar Mass Readout.
\end{proof}
\begin{definition}[speed In Meters Per Second]
  \label{def:physics:phyx-mini-0564:phyxminiproblems-problemphyxmini0564-speedinmeterspersecond}
  \lean{PhyXMiniProblems.ProblemPhyXMini0564.speedInMetersPerSecond}
  Read a physical speed in metres per second
\end{definition}
\begin{proof}
  This is the declared model definition of speed In Meters Per Second.
\end{proof}
\begin{definition}[rotation Frequency Readout]
  \label{def:physics:phyx-mini-0564:phyxminiproblems-problemphyxmini0564-rotationfrequencyreadout}
  \lean{PhyXMiniProblems.ProblemPhyXMini0564.rotationFrequencyReadout}
  \uses{def:physics:phyx-mini-0564:phyxminiproblems-problemphyxmini0564-rotationfrequencyquantity}
  Read drum rotation frequency in revolutions per selected time unit
\end{definition}
\begin{proof}
  This is the declared model definition of rotation Frequency Readout.
\end{proof}
\begin{definition}[molar Gas Constant In SI]
  \label{def:physics:phyx-mini-0564:phyxminiproblems-problemphyxmini0564-molargasconstantinsi}
  \lean{PhyXMiniProblems.ProblemPhyXMini0564.molarGasConstantInSI}
  \uses{def:physics:phyx-mini-0564:phyxminiproblems-problemphyxmini0564-molargasconstantquantity}
  Read the molar gas constant in coherent SI joules per mole-kelvin
\end{definition}
\begin{proof}
  This is the declared model definition of molar Gas Constant In SI.
\end{proof}
\begin{definition}[Molecular Species]
  \label{def:physics:phyx-mini-0564:phyxminiproblems-problemphyxmini0564-molecularspecies}
  \lean{PhyXMiniProblems.ProblemPhyXMini0564.MolecularSpecies}
  Molecular species admitted into the drum
\end{definition}
\begin{proof}
  This is the declared model definition of Molecular Species.
\end{proof}
\begin{definition}[Figure Label]
  \label{def:physics:phyx-mini-0564:phyxminiproblems-problemphyxmini0564-figurelabel}
  \lean{PhyXMiniProblems.ProblemPhyXMini0564.FigureLabel}
  Literal apparatus and marker labels visible in the supplied bitmap
\end{definition}
\begin{proof}
  This is the declared model definition of Figure Label.
\end{proof}
\begin{definition}[Plate Marker]
  \label{def:physics:phyx-mini-0564:phyxminiproblems-problemphyxmini0564-platemarker}
  \lean{PhyXMiniProblems.ProblemPhyXMini0564.PlateMarker}
  The two marked endpoints relevant to position along the curved plate
\end{definition}
\begin{proof}
  This is the declared model definition of Plate Marker.
\end{proof}
\begin{definition}[Plate Traversal]
  \label{def:physics:phyx-mini-0564:phyxminiproblems-problemphyxmini0564-platetraversal}
  \lean{PhyXMiniProblems.ProblemPhyXMini0564.PlateTraversal}
  Direction of traversal along the visible front portion of the plate
\end{definition}
\begin{proof}
  This is the declared model definition of Plate Traversal.
\end{proof}
\begin{definition}[Stern Rotating Drum Figure]
  \label{def:physics:phyx-mini-0564:phyxminiproblems-problemphyxmini0564-sternrotatingdrumfigure}
  \lean{PhyXMiniProblems.ProblemPhyXMini0564.SternRotatingDrumFigure}
  \uses{def:physics:phyx-mini-0564:phyxminiproblems-problemphyxmini0564-figurelabel, def:physics:phyx-mini-0564:phyxminiproblems-problemphyxmini0564-platetraversal}
  Qualitative geometry transcribed from the primary image. The image contains no numerical impact offset or answer value
\end{definition}
\begin{proof}
  This is the declared model definition of Stern Rotating Drum Figure.
\end{proof}
\begin{definition}[Stern Speed Distribution Setup]
  \label{def:physics:phyx-mini-0564:phyxminiproblems-problemphyxmini0564-sternspeeddistributionsetup}
  \lean{PhyXMiniProblems.ProblemPhyXMini0564.SternSpeedDistributionSetup}
  \uses{def:physics:phyx-mini-0564:phyxminiproblems-problemphyxmini0564-lengthquantity, def:physics:phyx-mini-0564:phyxminiproblems-problemphyxmini0564-timequantity, def:physics:phyx-mini-0564:phyxminiproblems-problemphyxmini0564-temperaturequantity, def:physics:phyx-mini-0564:phyxminiproblems-problemphyxmini0564-molarmassquantity, def:physics:phyx-mini-0564:phyxminiproblems-problemphyxmini0564-rotationfrequencyquantity, def:physics:phyx-mini-0564:phyxminiproblems-problemphyxmini0564-molargasconstantquantity, def:physics:phyx-mini-0564:phyxminiproblems-problemphyxmini0564-molecularspecies, def:physics:phyx-mini-0564:phyxminiproblems-problemphyxmini0564-sternrotatingdrumfigure}
  Physical quantities and observables of the experiment. The speed-to-impact map and the mean speed are independent fields. They are related to the input data only by the governing-law interfaces below, never by an answer-valued definition
\end{definition}
\begin{proof}
  This is the declared model definition of Stern Speed Distribution Setup.
\end{proof}
\begin{definition}[mean Speed Flight Time]
  \label{def:physics:phyx-mini-0564:phyxminiproblems-problemphyxmini0564-meanspeedflighttime}
  \lean{PhyXMiniProblems.ProblemPhyXMini0564.meanSpeedFlightTime}
  \uses{def:physics:phyx-mini-0564:phyxminiproblems-problemphyxmini0564-timequantity, def:physics:phyx-mini-0564:phyxminiproblems-problemphyxmini0564-sternspeeddistributionsetup}
  Flight time of molecules traveling at the Maxwell mean speed
\end{definition}
\begin{proof}
  This is the declared model definition of mean Speed Flight Time.
\end{proof}
\begin{definition}[mean Speed Strike Offset Revolutions]
  \label{def:physics:phyx-mini-0564:phyxminiproblems-problemphyxmini0564-meanspeedstrikeoffsetrevolutions}
  \lean{PhyXMiniProblems.ProblemPhyXMini0564.meanSpeedStrikeOffsetRevolutions}
  \uses{def:physics:phyx-mini-0564:phyxminiproblems-problemphyxmini0564-sternspeeddistributionsetup}
  Dimensionless plate offset from A, measured as a fraction of one turn
\end{definition}
\begin{proof}
  This is the declared model definition of mean Speed Strike Offset Revolutions.
\end{proof}
\begin{definition}[mean Speed Strike Arc Distance]
  \label{def:physics:phyx-mini-0564:phyxminiproblems-problemphyxmini0564-meanspeedstrikearcdistance}
  \lean{PhyXMiniProblems.ProblemPhyXMini0564.meanSpeedStrikeArcDistance}
  \uses{def:physics:phyx-mini-0564:phyxminiproblems-problemphyxmini0564-lengthquantity, def:physics:phyx-mini-0564:phyxminiproblems-problemphyxmini0564-sternspeeddistributionsetup}
  Physical arc distance along the plate from A for mean-speed molecules
\end{definition}
\begin{proof}
  This is the declared model definition of mean Speed Strike Arc Distance.
\end{proof}
\begin{definition}[Matches Stern Experiment Scenario]
  \label{def:physics:phyx-mini-0564:phyxminiproblems-problemphyxmini0564-matchessternexperimentscenario}
  \lean{PhyXMiniProblems.ProblemPhyXMini0564.MatchesSternExperimentScenario}
  \uses{def:physics:phyx-mini-0564:phyxminiproblems-problemphyxmini0564-sternspeeddistributionsetup}
  The qualitative experimental scenario stated in the prose
\end{definition}
\begin{proof}
  This is the declared model definition of Matches Stern Experiment Scenario.
\end{proof}
\begin{definition}[Matches Stern Problem Readouts]
  \label{def:physics:phyx-mini-0564:phyxminiproblems-problemphyxmini0564-matchessternproblemreadouts}
  \lean{PhyXMiniProblems.ProblemPhyXMini0564.MatchesSternProblemReadouts}
  \uses{def:physics:phyx-mini-0564:phyxminiproblems-problemphyxmini0564-lengthreadout, def:physics:phyx-mini-0564:phyxminiproblems-problemphyxmini0564-temperaturereadout, def:physics:phyx-mini-0564:phyxminiproblems-problemphyxmini0564-rotationfrequencyreadout, def:physics:phyx-mini-0564:phyxminiproblems-problemphyxmini0564-sternspeeddistributionsetup}
  Numerical readouts stated by the problem. No flight time, impact offset, arc distance, or displayed answer occurs here
\end{definition}
\begin{proof}
  This is the declared model definition of Matches Stern Problem Readouts.
\end{proof}
\begin{definition}[Uses Bismuth Dimer Reference Data]
  \label{def:physics:phyx-mini-0564:phyxminiproblems-problemphyxmini0564-usesbismuthdimerreferencedata}
  \lean{PhyXMiniProblems.ProblemPhyXMini0564.UsesBismuthDimerReferenceData}
  \uses{def:physics:phyx-mini-0564:phyxminiproblems-problemphyxmini0564-molarmassreadout, def:physics:phyx-mini-0564:phyxminiproblems-problemphyxmini0564-molargasconstantinsi, def:physics:phyx-mini-0564:phyxminiproblems-problemphyxmini0564-sternspeeddistributionsetup}
  Standard rounded classroom data for Bi₂ and the universal molar gas constant. These determine the Maxwell speed but do not state an impact position
\end{definition}
\begin{proof}
  This is the declared model definition of Uses Bismuth Dimer Reference Data.
\end{proof}
\begin{definition}[Matches Primary Stern Figure]
  \label{def:physics:phyx-mini-0564:phyxminiproblems-problemphyxmini0564-matchesprimarysternfigure}
  \lean{PhyXMiniProblems.ProblemPhyXMini0564.MatchesPrimarySternFigure}
  \uses{def:physics:phyx-mini-0564:phyxminiproblems-problemphyxmini0564-sternspeeddistributionsetup}
  Evidence read directly from the primary bitmap
\end{definition}
\begin{proof}
  This is the declared model definition of Matches Primary Stern Figure.
\end{proof}
\begin{definition}[Has Physical Stern Parameters]
  \label{def:physics:phyx-mini-0564:phyxminiproblems-problemphyxmini0564-hasphysicalsternparameters}
  \lean{PhyXMiniProblems.ProblemPhyXMini0564.HasPhysicalSternParameters}
  \uses{def:physics:phyx-mini-0564:phyxminiproblems-problemphyxmini0564-lengthreadout, def:physics:phyx-mini-0564:phyxminiproblems-problemphyxmini0564-timereadout, def:physics:phyx-mini-0564:phyxminiproblems-problemphyxmini0564-temperaturereadout, def:physics:phyx-mini-0564:phyxminiproblems-problemphyxmini0564-molarmassreadout, def:physics:phyx-mini-0564:phyxminiproblems-problemphyxmini0564-speedinmeterspersecond, def:physics:phyx-mini-0564:phyxminiproblems-problemphyxmini0564-rotationfrequencyreadout, def:physics:phyx-mini-0564:phyxminiproblems-problemphyxmini0564-molargasconstantinsi, def:physics:phyx-mini-0564:phyxminiproblems-problemphyxmini0564-sternspeeddistributionsetup, def:physics:phyx-mini-0564:phyxminiproblems-problemphyxmini0564-meanspeedflighttime}
  Positivity conditions selecting the intended physical regime
\end{definition}
\begin{proof}
  This is the declared model definition of Has Physical Stern Parameters.
\end{proof}
\begin{definition}[Satisfies Maxwell Mean Speed Law]
  \label{def:physics:phyx-mini-0564:phyxminiproblems-problemphyxmini0564-satisfiesmaxwellmeanspeedlaw}
  \lean{PhyXMiniProblems.ProblemPhyXMini0564.SatisfiesMaxwellMeanSpeedLaw}
  \uses{def:physics:phyx-mini-0564:phyxminiproblems-problemphyxmini0564-temperaturereadout, def:physics:phyx-mini-0564:phyxminiproblems-problemphyxmini0564-molarmassreadout, def:physics:phyx-mini-0564:phyxminiproblems-problemphyxmini0564-speedinmeterspersecond, def:physics:phyx-mini-0564:phyxminiproblems-problemphyxmini0564-molargasconstantinsi, def:physics:phyx-mini-0564:phyxminiproblems-problemphyxmini0564-sternspeeddistributionsetup}
  Maxwell's theoretical mean-speed relation ⟨v⟩ = sqrt (8 R T / (π M)) for a gas with molar mass M
\end{definition}
\begin{proof}
  This is the declared model definition of Satisfies Maxwell Mean Speed Law.
\end{proof}
\begin{definition}[Satisfies Rotating Drum Kinematics]
  \label{def:physics:phyx-mini-0564:phyxminiproblems-problemphyxmini0564-satisfiesrotatingdrumkinematics}
  \lean{PhyXMiniProblems.ProblemPhyXMini0564.SatisfiesRotatingDrumKinematics}
  \uses{def:physics:phyx-mini-0564:phyxminiproblems-problemphyxmini0564-lengthreadout, def:physics:phyx-mini-0564:phyxminiproblems-problemphyxmini0564-timereadout, def:physics:phyx-mini-0564:phyxminiproblems-problemphyxmini0564-speedinmeterspersecond, def:physics:phyx-mini-0564:phyxminiproblems-problemphyxmini0564-rotationfrequencyreadout, def:physics:phyx-mini-0564:phyxminiproblems-problemphyxmini0564-sternspeeddistributionsetup}
  Straight-line flight, uniform rotation, and circular plate geometry. These relations hold generically for every positive molecular speed; in particular, they do not contain 0.05021, choice D, or any problem-specific offset
\end{definition}
\begin{proof}
  This is the declared model definition of Satisfies Rotating Drum Kinematics.
\end{proof}
\begin{definition}[Satisfies Inverse Speed Sorting]
  \label{def:physics:phyx-mini-0564:phyxminiproblems-problemphyxmini0564-satisfiesinversespeedsorting}
  \lean{PhyXMiniProblems.ProblemPhyXMini0564.SatisfiesInverseSpeedSorting}
  \uses{def:physics:phyx-mini-0564:phyxminiproblems-problemphyxmini0564-speedinmeterspersecond, def:physics:phyx-mini-0564:phyxminiproblems-problemphyxmini0564-sternspeeddistributionsetup}
  Faster molecules have a smaller rotation-normalized offset from A
\end{definition}
\begin{proof}
  This is the declared model definition of Satisfies Inverse Speed Sorting.
\end{proof}
\begin{definition}[Satisfies Densitometer Law]
  \label{def:physics:phyx-mini-0564:phyxminiproblems-problemphyxmini0564-satisfiesdensitometerlaw}
  \lean{PhyXMiniProblems.ProblemPhyXMini0564.SatisfiesDensitometerLaw}
  \uses{def:physics:phyx-mini-0564:phyxminiproblems-problemphyxmini0564-sternspeeddistributionsetup}
  Deposit density is proportional to molecule count at each plate offset
\end{definition}
\begin{proof}
  This is the declared model definition of Satisfies Densitometer Law.
\end{proof}
\begin{definition}[Answer Choice]
  \label{def:physics:phyx-mini-0564:phyxminiproblems-problemphyxmini0564-answerchoice}
  \lean{PhyXMiniProblems.ProblemPhyXMini0564.AnswerChoice}
  Labels of the four revolution-fraction choices in the source
\end{definition}
\begin{proof}
  This is the declared model definition of Answer Choice.
\end{proof}
\begin{definition}[displayed Offset Revolutions]
  \label{def:physics:phyx-mini-0564:phyxminiproblems-problemphyxmini0564-displayedoffsetrevolutions}
  \lean{PhyXMiniProblems.ProblemPhyXMini0564.displayedOffsetRevolutions}
  \uses{def:physics:phyx-mini-0564:phyxminiproblems-problemphyxmini0564-answerchoice}
  Displayed normalized plate distances from A, in revolutions
\end{definition}
\begin{proof}
  This is the declared model definition of displayed Offset Revolutions.
\end{proof}
\begin{definition}[recorded Dataset Answer]
  \label{def:physics:phyx-mini-0564:phyxminiproblems-problemphyxmini0564-recordeddatasetanswer}
  \lean{PhyXMiniProblems.ProblemPhyXMini0564.recordedDatasetAnswer}
  \uses{def:physics:phyx-mini-0564:phyxminiproblems-problemphyxmini0564-answerchoice}
  The answer label recorded by the source dataset
\end{definition}
\begin{proof}
  This is the declared model definition of recorded Dataset Answer.
\end{proof}
\begin{definition}[Agrees Within Revolutions]
  \label{def:physics:phyx-mini-0564:phyxminiproblems-problemphyxmini0564-agreeswithinrevolutions}
  \lean{PhyXMiniProblems.ProblemPhyXMini0564.AgreesWithinRevolutions}
  A dimensionless modeled offset agrees with a display within tolerance
\end{definition}
\begin{proof}
  This is the declared model definition of Agrees Within Revolutions.
\end{proof}
\begin{definition}[Is Unique Nearest Displayed Offset]
  \label{def:physics:phyx-mini-0564:phyxminiproblems-problemphyxmini0564-isuniquenearestdisplayedoffset}
  \lean{PhyXMiniProblems.ProblemPhyXMini0564.IsUniqueNearestDisplayedOffset}
  \uses{def:physics:phyx-mini-0564:phyxminiproblems-problemphyxmini0564-sternspeeddistributionsetup, def:physics:phyx-mini-0564:phyxminiproblems-problemphyxmini0564-meanspeedstrikeoffsetrevolutions, def:physics:phyx-mini-0564:phyxminiproblems-problemphyxmini0564-answerchoice, def:physics:phyx-mini-0564:phyxminiproblems-problemphyxmini0564-displayedoffsetrevolutions}
  A choice is strictly nearer to the modeled offset than every other choice
\end{definition}
\begin{proof}
  This is the declared model definition of Is Unique Nearest Displayed Offset.
\end{proof}
\begin{lemma}[mean Speed Strike Offset eq rotation Rate mul distance div speed]
  \label{lem:physics:phyx-mini-0564:phyxminiproblems-problemphyxmini0564-meanspeedstrikeoffset-eq-rotationrate-mul-distance-div-speed}
  \lean{PhyXMiniProblems.ProblemPhyXMini0564.meanSpeedStrikeOffset_eq_rotationRate_mul_distance_div_speed}
  \uses{def:physics:phyx-mini-0564:phyxminiproblems-problemphyxmini0564-lengthreadout, def:physics:phyx-mini-0564:phyxminiproblems-problemphyxmini0564-speedinmeterspersecond, def:physics:phyx-mini-0564:phyxminiproblems-problemphyxmini0564-rotationfrequencyreadout, def:physics:phyx-mini-0564:phyxminiproblems-problemphyxmini0564-sternspeeddistributionsetup, def:physics:phyx-mini-0564:phyxminiproblems-problemphyxmini0564-meanspeedstrikeoffsetrevolutions, def:physics:phyx-mini-0564:phyxminiproblems-problemphyxmini0564-hasphysicalsternparameters, def:physics:phyx-mini-0564:phyxminiproblems-problemphyxmini0564-satisfiesrotatingdrumkinematics}
  Combining straight-line flight with uniform rotation gives the generic time-of-flight expression for the normalized impact offset
\end{lemma}
\begin{proof}
  The conclusion follows from the governing relations and figure readouts named in the statement of mean Speed Strike Offset eq rotation Rate mul distance div speed.
\end{proof}
% --- Archon declaration topology end ---
% --- Archon physics formalization source end ---
